\documentclass[12pt,a4paper]{article}

% --- Codificación y bibliografía del documento ---
\usepackage[utf8]{inputenc}
\usepackage[T1]{fontenc}
\usepackage[spanish]{babel}

% --- Gráficos y enlaces ---
\usepackage{graphicx}
\usepackage{hyperref}

% --- Fragmentos de código ---
\usepackage{listings}
\usepackage{xcolor}
\lstset{
    basicstyle=\ttfamily\small,
    breaklines=true,
    frame=single,
    numbers=left,
    numbersep=5pt,
    commentstyle=\color{gray},
    keywordstyle=\color{blue},
    stringstyle=\color{teal}
}

% --- Bibliografía (natbib + bibtex, referencias.bib) ---
\usepackage{natbib}

\hypersetup{
    colorlinks=true,
    linkcolor=black,
    urlcolor=blue,
    citecolor=black
}

\begin{document}

% ============================================================
% PORTADA
% ============================================================
\begin{titlepage}
    \centering

    \vspace*{1cm}
    {\large Universidad Técnica Estatal de Quevedo\par}
    {\large Facultad de Ciencias de la Computación y Diseño Digital\par}
    {\large Carrera de Ingeniería de Software\par}

    \vspace{2.5cm}
    {\Huge\bfseries Consumo de Web Services REST desde un Front-end AlpineJS --- Proyecto ArtiSync\par}

    \vspace{2cm}
    {\large Asignatura: Aplicaciones Web (5to nivel) --- Segundo Corte\par}

    \vspace{1.5cm}
    {\large\bfseries Integrantes:\par}
    {\large Figueroa Morales Bryan Javier\par}
    {\large Loor Medranda Marlon Taylor\par}

    \vspace{1.5cm}
    {\large\bfseries Docente:\par}
    {\large Dr. Gleiston Cicerón Guerrero Ulloa, Ph.D.\par}

    \vspace{1.5cm}
    {\large Julio 2026\par}

    \vspace{1.5cm}
    {\small Repositorio: \url{https://github.com/011011bjf/GA-Exposici-n-Segundo-Corte-Frameworks-del-Back-End-}\par}

    \vfill
\end{titlepage}

% ============================================================
% CONTENIDO
% ============================================================

\section{Introducción}
% TODO: redactar introducción (contexto del proyecto ArtiSync, objetivo del informe, alcance del consumo REST).

\section{Paso a paso del consumo de la API REST}
% TODO: documentar, con capturas y fragmentos de código (login, registro, 2FA, CRUD de usuarios), cómo el frontend AlpineJS consume cada endpoint del backend.

\section{Conclusiones}
% TODO: redactar conclusiones del proyecto.

% ============================================================
% REFERENCIAS
% ============================================================
% TODO: agregar \cite{} en el cuerpo del documento a medida que se redacta;
% \nocite{*} se usa temporalmente para verificar que las 9 entradas del .bib compilan.
\nocite{*}
\bibliographystyle{plain}
\bibliography{referencias}

\end{document}
